\chapter{CONCLUSIONS AND FUTURE RESEARCH DIRECTIONS}

\section{Summary of main findings}

This thesis investigates user-dissatisfaction-based static bike repositioning, with a focus on service-quality-aware decomposition and the integrated handling of broken bikes through truck collection and on-site repair. The research progresses from a UDF-aware decomposition framework for the conventional multi-vehicle static bike repositioning problem in Chapter 3, to an integrated broken-bike and on-site-repair model in Chapter 4, and finally to an EUDF-aware large-scale decomposition framework for the extended problem in Chapter 5.

Chapter 3 develops a user-dissatisfaction-aware cluster-first, route-second decomposition framework for the multi-vehicle static bike repositioning problem. The clustering model jointly accounts for estimated UDF reduction and per-vehicle routing-time feasibility. The Activated Transfer Selection procedure efficiently evaluates candidate clusters by selecting time-feasible bike transfers, while the Hybrid Genetic Search for UDF-aware Clustering (HGS-UC) solves the resulting station-partitioning problem at city scale. The computational experiments show that HGS-UC obtains near-optimal solutions on 50-station instances, produces better feasible solutions than the solver incumbents on most 100-station tests, and achieves statistically significant improvements over five baseline clustering methods on networks with 200--1,000 stations.

Chapter 4 introduces a static bike repositioning problem with broken bikes and on-site repair. The problem is formulated as a time-indexed mixed-integer linear programming model that jointly determines truck routes, repairer routes, and station-level quantities for repairing, loading, and unloading bikes. Its objective combines the penalty cost of EUDF-based user dissatisfaction with the cost of truck CO$_2$ emissions within a common operational time budget. The proposed HGSADC-SBC heuristic combines Hybrid Genetic Search with Adaptive Diversity Control for route construction and a station-budget-constrained heuristic for loading, unloading, and repair decisions. The computational results show that the heuristic matches the optimal solutions obtained by Gurobi on small instances of up to 15 stations and produces better feasible solutions on networks with 60--500 stations within substantially less computational time. The experiments further show that the cost-effectiveness of introducing on-site repairers decreases when repair times are long or the proportion of broken bikes is low.

Chapter 5 extends the service-quality-aware decomposition approach to the static bike repositioning problem with broken bikes and on-site repair. The Extended Activated Transfer Selection procedure estimates the reduction in EUDF associated with usable-bike transfers, broken-bike collection, and on-site repair within each candidate cluster. HGS-EUC then uses these estimates to search for station partitions, after which HGSADC-SBC is applied to the cluster-level routing and quantity-planning problems. Experiments on instances with up to 2,000 stations show that the proposed framework obtains lower feasible objective values than the Gurobi upper bounds on the small benchmark instances and outperforms the monolithic HGSADC-SBC heuristic on all tested small and large instances. Its advantage also remains positive across different truck capacities, operational time budgets, and broken-bike rates.

\section{Contributions of the thesis}

This thesis makes three main contributions. First, it establishes a service-quality-aware decomposition approach for static bike repositioning. Rather than treating clustering as a purely geographic or demand-based preprocessing step, the proposed frameworks link station partitioning directly to the user-dissatisfaction objective of the subsequent repositioning problem. Chapter 3 develops this principle for the one-dimensional UDF setting, while Chapter 5 extends it to the bivariate EUDF setting and accounts for the coupled effects of usable-bike transfers, broken-bike collection, and on-site repair.

Second, the thesis develops an integrated operational model for static bike repositioning with broken bikes and on-site repair. The formulation jointly represents usable-bike repositioning, broken-bike collection, truck routing, repairer routing, and station-level repair decisions within the same planning horizon. It evaluates service outcomes through the EUDF while accounting for the emissions generated by truck operations. The associated computational analysis also identifies the conditions under which introducing on-site repairers is operationally cost-effective.

Third, the thesis develops a layered computational methodology for solving large-scale problems. ATS and EATS provide efficient surrogate evaluations of candidate clusters without requiring a detailed routing problem to be solved for every candidate partition. HGS-UC and HGS-EUC search for service-quality-aware station partitions, while HGSADC-SBC addresses the detailed routing and station-level quantity decisions. Together, these methods connect exact benchmark formulations with scalable heuristic planning and enable the studied problems to be addressed on networks containing up to 2,000 stations.

\section{Future research directions}

Several directions can extend this research. First, the deterministic static models can be extended to dynamic and stochastic settings. Incorporating uncertain rental and return demand, uncertain discovery of broken bikes, and real-time station inventories would support rolling-horizon or intra-day planning. Such an extension would require the UDF or EUDF values, station partitions, and vehicle and repairer routes to be updated as new information becomes available.

Second, future models can represent repair operations and bike movements in greater operational detail. Relevant extensions include heterogeneous fault types, repairer skills, uncertain repair durations, repairer shift scheduling, explicit depot-repair capacities and completion times, and preventive-maintenance decisions. Continuous-time formulations could reduce the approximation introduced by time discretization, while more flexible quantity-planning methods could allow temporary station inventories, repeated station visits, or intermediate bike transfers.

Third, coordination across clusters can be strengthened. The frameworks in Chapters 3 and 5 fix the station partition before detailed routing and quantity planning. Adaptive repartitioning, boundary-station exchanges, and limited cross-cluster operations could retain the computational benefits of decomposition while recovering beneficial actions excluded by a fixed partition. Promising candidate partitions could also be re-evaluated through stronger routing subproblem approximations to improve consistency between the clustering and routing stages.

Finally, the service-quality-aware decomposition frameworks in Chapters 3 and 5 can be extended to multi-criteria objectives. In addition to user dissatisfaction, cluster formation and cluster-level planning may jointly consider truck emissions, workload balance among vehicles and repairers, fairness across stations, and robustness to operational disruptions. Incorporating these criteria into the cluster-evaluation and search procedures would broaden the operational applicability of the proposed decomposition methods.